% ==============================================================================
% BAB IV. PENUTUP (BERWARNA RESMI BPS)
% ==============================================================================

\chapter{PENUTUP}

\section{Gambaran Kendala dan Solusi}

\begin{enumerate}[label=\textbf{\textcolor{bpsnavy}{\arabic*.}}]
  \item \textbf{Keterbatasan Sumber Daya Manusia (SDM) dan Literasi Data}
  \begin{itemize}
    \item \textbf{\textcolor{ribbonred}{Kendala:}}
    \begin{itemize}
      \item Masih terbatasnya pemahaman teknis aparatur desa dan kader terkait istilah-istilah statistik serta metodologi pengumpulan data.
      \item Adanya perbedaan persepsi dan ketidakpahaman dari beberapa Kepala Jaga/Meweteng saat mengumpulkan data di tingkat lingkungan.
    \end{itemize}
    \item \textbf{\textcolor{bpsgreen}{Solusi:}}
    \begin{itemize}
      \item Menyelenggarakan pelatihan dan bimbingan teknis (\textit{coaching}) secara bertahap bagi Agen Statistik Desa dan aparatur desa.
      \item Pembina Desa Cantik dari BPS bersama Agen Statistik melakukan pertemuan pendampingan langsung untuk memberikan penjelasan mengenai definisi operasional dan istilah statistik.
    \end{itemize}
  \end{itemize}

  \item \textbf{Ketersediaan dan Pengelolaan Data Desa}
  \begin{itemize}
    \item \textbf{\textcolor{ribbonred}{Kendala:}}
    \begin{itemize}
      \item Kelengkapan data variabel desa masih belum seragam, di mana sebagian data sektor publik sudah tersedia namun sebagian lainnya belum terdokumentasi dengan baik.
      \item Sistem penyimpanan data sebelumnya masih bersifat parsial/manual sehingga berisiko hilang atau sulit diakses kembali saat dibutuhkan.
    \end{itemize}
    \item \textbf{\textcolor{bpsgreen}{Solusi:}}
    \begin{itemize}
      \item Menerapkan mekanisme pengolahan dan tabulasi data secara terstruktur menggunakan aplikasi spreadsheet sederhana atau sistem pencatatan terpusat.
      \item Mengoptimalkan inventarisasi data profil desa agar selaras dengan prinsip Satu Data Indonesia (SDI).
    \end{itemize}
  \end{itemize}

  \item \textbf{Infrastruktur Teknologi Informasi}
  \begin{itemize}
    \item \textbf{\textcolor{ribbonred}{Kendala:}}
    \begin{itemize}
      \item Ketersediaan sarana pendukung (seperti komputer dan jaringan internet) di tingkat desa belum sepenuhnya optimal untuk mendukung pengolahan data digital secara cepat.
    \end{itemize}
    \item \textbf{\textcolor{bpsgreen}{Solusi:}}
    \begin{itemize}
      \item Menggabungkan metode pendataan secara \textit{offline} (manual/kertas) terlebih dahulu, kemudian diolah dan diunggah secara berkala ketika akses sistem/perangkat tersedia.
      \item Mendorong pemanfaatan fasilitas internet kantor desa secara efisien untuk pembuatan media penyajian data digital (website/infografis desa).
    \end{itemize}
  \end{itemize}

  \item \textbf{Kesadaran dan Partisipasi Masyarakat}
  \begin{itemize}
    \item \textbf{\textcolor{ribbonred}{Kendala:}}
    \begin{itemize}
      \item Sebagian masyarakat belum menyadari penuh pentingnya akurasi data statistik, sehingga partisipasi saat pendataan pendukung terkadang kurang maksimal.
    \end{itemize}
    \item \textbf{\textcolor{bpsgreen}{Solusi:}}
    \begin{itemize}
      \item Mengintegrasikan sosialisasi literasi data ke dalam pertemuan warga atau kegiatan resmi desa.
      \item Melibatkan tokoh masyarakat dan perangkat desa terdepan untuk memberikan pemahaman bahwa data akurat akan bermuara pada program pembangunan desa yang lebih tepat sasaran.
    \end{itemize}
  \end{itemize}
\end{enumerate}

\begin{table}[htbp]
\centering
\small
\caption{Matriks Kendala dan Solusi}
\label{tab:matriks_kendala}
\vspace{0.2cm}
\begin{tabularx}{\textwidth}{|c|Y|Y|Y|}
\hline
\rowcolor{bpsnavy}
\textcolor{white}{\textbf{No}} & \centering\textcolor{white}{\textbf{Bidang Kendala}} & \centering\textcolor{white}{\textbf{Permasalahan Utama}} & \centering\textcolor{white}{\textbf{Solusi / Tindak Lanjut}} \tabularnewline
\hline
\rowcolor{white}
1 & SDM \& Pemahaman & Kurang memahami istilah teknis statistik. & Pendampingan intensif dari Pembina BPS \& Agen Desa. \tabularnewline
\hline
\rowcolor{bpslight!50}
2 & Kelengkapan Data & Sebagian data belum tersedia/belum terintegrasi. & Inventarisasi data profil desa \& standarisasi format SDI. \tabularnewline
\hline
\rowcolor{white}
3 & Infrastruktur / IT & Sarana digital \& konektivitas belum memadai. & Pengolahan bertahap (\textit{offline-to-online}) \& optimalisasi sarana desa. \tabularnewline
\hline
\rowcolor{bpslight!50}
4 & Pengelolaan Data & Data belum terstruktur rapi. & Penerapan manajemen data \& pembuatan publikasi infografis. \tabularnewline
\hline
\end{tabularx}
\end{table}

\section{Kesimpulan dan Saran}

\subsection*{Kesimpulan:}

Kegiatan Desa Cinta Statistik (Desa Cantik) merupakan program strategis dari BPS untuk meningkatkan kapasitas desa dalam mengelola dan memanfaatkan data statistik. Tujuannya adalah agar desa-desa memiliki kemampuan mengumpulkan, mengelola, dan menggunakan data secara mandiri sebagai dasar dalam perencanaan dan pengambilan kebijakan yang lebih efektif. Meskipun program ini dihadapkan pada berbagai kendala seperti keterbatasan sumber daya manusia, infrastruktur teknologi, partisipasi masyarakat, dan pengelolaan data, solusi yang tepat dapat mengatasi tantangan tersebut. Dengan pelatihan, peningkatan akses teknologi, serta sosialisasi yang efektif, desa-desa diharapkan dapat mencapai kemandirian dalam pengelolaan data dan memanfaatkannya untuk pembangunan yang lebih baik.

\subsection*{Saran:}

\begin{enumerate}[label=\textbf{\textcolor{bpsnavy}{\arabic*.}}]
  \item \textbf{Perluasan Jangkauan Pelatihan:} BPS sebaiknya memperluas cakupan pelatihan statistik hingga ke seluruh perangkat desa dan melibatkan relawan lokal. Penggunaan modul yang mudah dipahami akan sangat membantu meningkatkan pemahaman.
  \item \textbf{Penguatan Infrastruktur Teknologi:} Pemerintah dan pihak terkait perlu bekerja sama untuk meningkatkan akses teknologi di desa, seperti menyediakan komputer dan akses internet yang memadai, terutama di desa terpencil.
  \item \textbf{Sosialisasi yang Lebih Intensif:} Sosialisasi tentang pentingnya data statistik harus dilakukan secara lebih intensif dengan melibatkan tokoh masyarakat dan media lokal untuk mendorong partisipasi aktif warga dalam pendataan.
  \item \textbf{Monitoring dan Evaluasi Berkelanjutan:} Implementasi program Desa Cantik perlu didukung oleh monitoring dan evaluasi secara berkala untuk memastikan pengelolaan data berjalan dengan baik dan sesuai tujuan awal, serta mengidentifikasi area yang perlu ditingkatkan.
\end{enumerate}
